%! Boxplots and Comparing Distributions — Unit 1, Lesson 1.6 — Lesson Plan
% GRADUAL-RELEASE plan (warm-up / guided notes and practice / debrief / close and start homework).
% REGENERATED from the retired EFFL shape (LESSON_SHAPE.md §7): experience/ and experience_key/
% are gone, notes/ + notes_key/ carry the whole release, and the spoiler rule is DEAD — use the
% formal vocabulary freely everywhere, students included.
% THERE IS NO GROUP ACTIVITY IN THIS COURSE — the guided notes carry the release row by row in
% one Main Ideas / Notes table: I do (each row's definition lines and first problem) -> we do
% (the rest of the grid, then the Guided Practice row). THERE IS NO INDEPENDENT PRACTICE SET —
% the homework is the individual practice, and the period ends with students starting it.
% Teacher-facing. Every box is `skillbox` (breakable) with a \boxguard ahead of it; this course
% has no `fixedskillbox`, no tier boxes, and no exit ticket.
% THIS LESSON MERGES CED TOPICS 1.8 AND 1.9 into one period, so the notes table carries both:
% rows 1, 2 and 4 build and read ONE boxplot (1.8), row 3 and Guided Practice compare TWO (1.9).
% Row 4, One Quarter, is the CRUX ROW and therefore the LAST INSTRUCTION ROW before Guided
% Practice (LESSON_SHAPE.md §1); the comparison row comes before it.
\documentclass[10pt]{article}
\usepackage{apstats-article}
\usepackage{apstats-boxes}

\newcommand{\UnitNumberName}{Unit 1: Exploring One-Variable Data \quad}
\newcommand{\LessonNumberName}{Lesson 1.6: Boxplots and Comparing Distributions}

\begin{document}
\begin{center}
    {\Large\bfseries \CourseName \\
    {\small\bfseries \UnitNumberName \LessonNumberName}}
\end{center}

\begin{tcolorbox}[colback=sky, colframe=navy, boxrule=0.9pt, arc=2mm,
    left=2mm, right=2mm, top=1.5mm, bottom=1.5mm]
    \textbf{Primary Objective:} Students will assemble the five-number summary for a quantitative
    variable, read a boxplot as four sections each holding a quarter of the data, and compare two
    distributions on center, variability, and unusual features --- in context, with explicit
    comparative language.
    \\[2pt]
    \textbf{CED alignment:} Topics 1.8 and 1.9 \quad$\cdot$\quad Big Idea: UNC (Uncertainty)
    \quad$\cdot$\quad Skills 2.A, 2.B, 2.D \quad$\cdot$\quad LO UNC-1.L, UNC-1.M, UNC-1.N,
    UNC-1.O \quad$\cdot$\quad EK UNC-1.L.1, 1.L.2, 1.M.1, 1.M.2, 1.N.1, 1.O.1
    \\[2pt]
    \textbf{Lesson model:} \emph{Gradual release --- I do, we do, you do --- all of it inside
    the guided notes.} There is no group activity. The notes name the vocabulary as they teach
    it and deliver the instruction row by row in one \emph{Main Ideas / Notes} table --- each
    idea a definition to complete and a grid of short problems, the first worked by you and the
    rest by the class; the last row, Guided Practice (\emph{Two Science Sections}), is worked
    together with students holding the pen; \textbf{the homework is the individual practice}, and
    students start it in class while you circulate. The whole-class debrief is \emph{spoken} ---
    there is no transfer set and no reflection box on the page.
    \\[2pt]
    \textbf{Note on the ``no sketch'' rule:} \textbf{students never draw a boxplot in this
    lesson.} Every display in the notes, the AP Practice page, and the homework is pre-drawn.
    Where a question says ``construct,'' the construction is filling in a table of five numbers
    or counting dots into stated sections --- never plotting an axis. Drawing boxplots by hand
    eats the period; the exam asks students to \emph{read} and \emph{compare} them.
    \\[2pt]
    \textbf{Note on the merge:} this period carries Topic 1.8 (build and read one boxplot) in
    notes rows 1, 2 and 4 and Topic 1.9 (compare two or more) in row 3 and Guided Practice. Protect
    the 34 minutes of notes; the extra comparison reps live in the AP Practice page and the
    homework, not in a sixth row.
\end{tcolorbox}

\renewcommand{\arraystretch}{1.4}
\boxguard
\begin{skillbox}[Learning Targets \& Key Understandings]{goldbox}
\begin{tabularx}{\linewidth}{@{}X|X@{}}
\textbf{Learning targets (student language)} & \textbf{Key understandings (the ``why'')} \\[2pt]
\hline
\vspace{-6pt}
\begin{itemize}[leftmargin=1.1em, nosep, after=\vspace{2pt}]
  \item I can report the \textbf{five-number summary} of a data set and point to each of the
        five numbers on a \textbf{boxplot}.
  \item I can say what a wide section of a boxplot does and does not tell me.
  \item I can decide with the $1.5\times$IQR rule whether a value is an outlier, and say how a
        \textbf{modified boxplot} draws it.
  \item I can compare two distributions from \textbf{parallel boxplots} on center, variability,
        and unusual features --- in context, with comparative language.
\end{itemize}
&
\vspace{-6pt}
\begin{itemize}[leftmargin=1.1em, nosep, after=\vspace{2pt}]
  \item The minimum, $Q_1$, the median, $Q_3$, and the maximum are the \textbf{five-number
        summary}; a boxplot is that summary drawn to scale, and nothing more (EK UNC-1.L.1,
        1.L.2).
  \item \textbf{Target misconception:} that a wide section of a boxplot holds more
        observations. \textbf{Every section holds the same quarter of the data.} Width reports
        spread, not count --- and this is the single most common boxplot error on the exam.
  \item \textbf{Second misconception:} that comparing two groups means saying which median is
        higher. The CED wants center, variability, \emph{and} unusual features, each with a
        comparative word and each in context (EK UNC-1.N.1, 1.O.1).
  \item Skew moves the mean off the median: right-skew pulls the mean right of it, left-skew
        pulls it left (EK UNC-1.M.2). And five numbers cost you everything between them --- a
        boxplot cannot show a gap, a cluster, or a second peak.
\end{itemize}
\end{tabularx}
\end{skillbox}

\boxguard
\begin{skillbox}[{Vocabulary, Concepts, \& Theorems --- taught directly in the guided notes}]{greenbox}
\begin{minipage}[t]{0.48\linewidth}
    \begin{itemize}
        \item \textbf{Five-number summary} --- minimum, $Q_1$, median, $Q_3$, maximum
              (EK UNC-1.L.1)
        \item \textbf{Boxplot} --- the five-number summary drawn to scale (EK UNC-1.L.2)
        \item \textbf{Box} --- spans $Q_1$ to $Q_3$: the middle 50\% of the data, with a line at
              the median (EK UNC-1.L.2)
        \item \textbf{Whisker} --- reaches from a quartile to the most extreme value that is
              \emph{not} an outlier (EK UNC-1.L.2)
    \end{itemize}
\end{minipage}
\hfill
\begin{minipage}[t]{0.48\linewidth}
    \begin{itemize}
        \item \textbf{Modified boxplot} --- outliers plotted with their own symbol beyond the
              whisker (EK UNC-1.L.2)
        \item \textbf{Parallel (side-by-side) boxplots} --- one plot per group on a shared axis
              (EK UNC-1.N.1)
        \item \textbf{Mean vs.\ median position} --- symmetric: close; right-skewed: mean to the
              right; left-skewed: mean to the left (EK UNC-1.M.2)
        \item \textbf{Carried forward} --- median, $Q_1$, $Q_3$, IQR, range, resistance, and the
              $1.5\times$IQR rule (1.5); shape and unusual features (1.4); dotplots and
              histograms (1.3)
    \end{itemize}
\end{minipage}
\end{skillbox}

% A tabularx never splits, so guard generously ahead of it.
\boxguard[30]
\begin{skillbox}[Lesson at a Glance --- \MeetingLength]{sky}
\renewcommand{\arraystretch}{1.25}
\begin{tabularx}{\linewidth}{@{}l c X X@{}}
\textbf{Phase} & \textbf{Min} & \textbf{Students} & \textbf{Teacher} \\
\hline
Warm-Up & 5 & Five numbers off an ordered list, a counting drill, and one $1.5\times$IQR check &
              Debrief item~1 aloud; put item~3's arithmetic on the board and leave it up all
              period \\
Guided Notes \& Practice & 34 & Fill the vocabulary box and the notes table: problems 1--14 row
              by row, then \emph{Two Science Sections} (15--20) together, pen in their hand &
              \textbf{I do} each row's definition lines and first problem, \textbf{we do} the
              rest of its grid (24) $\to$ \textbf{we do} Guided Practice (10) \\
Debrief & 8 & Pens down, papers up --- answer aloud, nothing to fill in & Open on an
              \emph{almost}-right answer to problem~11, then the count row of problem~10 and the
              comparison standard \\
Close \& Assign & 8 & \textbf{Start the homework in class}, alone and silent & Announce the paper
              homework --- due the first class after two study halls --- circulate for your
              formative read, preview Lesson~1.7 \\
\end{tabularx}
\end{skillbox}

\boxguard[20]
\begin{skillbox}[Warm-Up --- Activate Prior Knowledge (5 min)]{sky}
\begin{minipage}[t]{0.48\linewidth}
    \textbf{The three items and what each seeds}
    \begin{itemize}
        \item \textbf{1.} Nine ordered tree ages; report the least, $Q_1$, the median, $Q_3$, the
              greatest, plus the IQR and the range. \textbf{Seeds:} notes row~1 --- those same
              five numbers get their collective name, \emph{five-number summary}, in the first
              line of the table.
        \item \textbf{2.} Count twelve values into three stated stretches. \textbf{Seeds:}
              problem~10, the counting move the crux depends on, made mechanical first.
        \item \textbf{3.} One $1.5\times$IQR check with the arithmetic shown. \textbf{Seeds:}
              notes row~3's \texttt{work} block, the Route B rider at 58 minutes.
    \end{itemize}
\end{minipage}
\hfill
\begin{minipage}[t]{0.48\linewidth}
    \textbf{Running it}
    \begin{itemize}
        \item \textbf{Debrief item~1 aloud} and leave the five numbers on the board. You will
              point at them in the first minute of the notes and say ``these five have a name.''
        \item Item~2 is a drill --- 60 seconds. It exists so that counting is not the obstacle
              in problem~10, where the count \emph{is} the point.
        \item \textbf{Leave item~3's arithmetic on the board for the rest of the period.}
              Notes row~3 reuses it with different numbers, and a student who can look up the
              structure rather than re-derive it spends the time on the reasoning.
    \end{itemize}
\end{minipage}
\end{skillbox}

\boxguard
\begin{skillbox}[Guided Notes \& Practice --- I do, then we do (34 min)]{sky}
\textbf{This is the whole lesson.} There is no group activity, and there is no independent
practice set on this page: \textbf{the homework is the individual practice}, started in class.
The notes are one \emph{Main Ideas / Notes} table --- five rows, 20 short problems. \textbf{In
every row you fill the definition lines and work the first problem of the grid; the class works
the rest, pen in their hand.} The vocabulary box is filled \emph{as you go}, never in advance,
and the free-throw dotplot with its boxplot beneath it is the picture rows 1, 2 and~4 all
point back at.
\begin{multicols}{2}
\textbf{The Five Numbers (6 min).} Point at the warm-up's five numbers still on the board and
name them: \emph{five-number summary}. Fill the definition lines, write the quartile
\texttt{work} block, then the boxplot sentence --- ``five numbers, nothing else'' is the claim the
whole lesson rests on. Work problem~1 yourself (minimum 30, read off the leftmost dot); problems
2--3 are theirs. \textbf{Problem~2 is worth a hand count} --- which part of the boxplot marks
$Q_1$ --- because students who answer ``the line in the middle'' have not yet accepted that the
box has \emph{two} edges.

\textbf{Anatomy of a Boxplot (7 min).} Fill the box/whisker/outlier lines, then write the IQR and
fence \texttt{work} block on the free-throw totals: $10.5$ and $74.5$, and nothing outside them.
Work problem~4 yourself ($50.5 - 34.5 = 16$); 5--6 are theirs. \textbf{Problem~5 is the quiet one
that matters}: ten of the twenty sit inside the box, and a student who says ``most of them'' is
already showing you the misconception row~4 exists to break. Problem~6 must name the individual
values.

\textbf{Comparing Distributions (6 min).} Fill the checklist line --- center, variability,
unusual features, comparative word, in context --- and say plainly that describing one group then
the other earns nothing. Work problem~7 yourself (Route A's middle half is 5 minutes, Route B's
is 12); 8--9 are theirs. \textbf{The second \texttt{work} block is where the warm-up arithmetic
gets reused}: $32 + 1.5(12) = 50$, and $58 > 50$, so the whisker stops at 40.

\columnbreak
\textbf{One Quarter (8 min) --- the crux row, and the last one before Guided Practice.} Fill the
four-sections line and the 25\% line, then release \textbf{problem~10, the count row}. Do not let
anyone estimate: every boundary falls in a gap, so the four counts are exactly 5, 5, 5, 5. Write
them on the board with \textbf{25\% \quad 25\% \quad 25\% \quad 25\%} above them.
\textbf{Problem~11 is the trap} --- the classmate who says the wide part is where most of the
shots landed --- and \textbf{problem~12 is the counterweight}: width reports spread, never count.
Take a hand count on 11 before anyone speaks; pick one \emph{almost}-right answer off a desk for
the debrief. Problem~14 is the cost of the box, and it returns in the AP Practice page's MC~4.

\textbf{We do (about 10 min): Two Science Sections.} Problems 15--20, worked entirely together,
in a context the lesson has not used. Students hold the pen; you hold the questions: \emph{``You
wrote `Period 5 has an IQR of 16.' Compared with what?''} \quad \emph{``Which side of Period 2's
summary is stretched? So which way does it trail off?''} \textbf{Problem~20 is the AP-scored
form} --- one paragraph, in context, with comparative language --- and it is the one that cannot
be cut. If you only get 15, 16, and~20, that is the right trade.

Circulate and demand the reason, not the number:
\begin{itemize}[leftmargin=1.1em, nosep]
  \item \textbf{Problem~11, the crux:} ``Count them. How many in the wide section? How many in
        the narrow one? Now say that out loud.''
  \item \textbf{Problem~18:} ``Where would an outlier be drawn if there were one?''
  \item \textbf{Problem~20:} ``Cross out every word in your paragraph that is not comparing.''
\end{itemize}
While circulating, pick the papers for the debrief --- one \emph{almost}-right problem~11, and
one problem~20 that describes twice instead of comparing once. \textbf{Your formative read comes
twice}: here, and again in the last eight minutes as students start the homework.
\end{multicols}
\end{skillbox}

\boxguard
\begin{skillbox}[Debrief --- whole class, spoken (8 min)]{sky}
Pens down, papers up. \textbf{This debrief is spoken} --- there is no box at the end of the notes
to fill in, so it lives entirely on the board and in the room.
\begin{multicols}{2}
\begin{enumerate}[leftmargin=1.3em, nosep, itemsep=3pt]
  \item \textbf{Open on the \emph{almost}-right problem~11 you collected} --- typically ``the
        wide part is spread out, so it has fewer people in the middle.'' Half right, and the half
        that is wrong is the lesson. Put the count row from problem~10 beside it: \textbf{5 \quad
        5 \quad 5 \quad 5}, and above it \textbf{25\% \quad 25\% \quad 25\% \quad 25\%}.
  \item Say the sentence, write the sentence, and point at the fives while you do: \emph{``Every
        section of a boxplot holds the same quarter of the data. Width tells you how stretched out
        those values are --- it never tells you how many there are.''}
  \item Take problem~5 next: ten of twenty inside the box, ten outside. The box is not ``where
        the data is''; it is where the middle half is.
  \item Put up the two problem~20 paragraphs and name the difference out loud: one wrote two
        descriptions, one wrote a comparison. \textbf{Only the second scores}, however accurate
        the first is.
  \item Close on the cost, off problem~14: gaps, clusters, and a second peak all vanish inside
        the box. If someone asks whether there are two kinds of night, ask for a dotplot.
\end{enumerate}
\columnbreak
\textbf{Two things to demand aloud.} First, \textbf{make a student state the quarter rule without
using the word ``wide''} --- that constraint is the point, because ``the wide part has more'' is
the sentence they walked in with. Second, \textbf{make them produce one comparison sentence about
a context not on any page today}, with a comparative word and a unit in it; that is where transfer
shows up now that there is no transfer set on paper.

\textbf{Connect back to Lesson~1.4 if time allows:} skew was a shape word there and it is a
position rule here --- right-skew puts the mean right of the median (problem~9). One line is
enough; it returns in 1.7.

\textbf{If the debrief runs short}, cut items~3 and~5. \textbf{Item~1 is the one that cannot be
cut}, and do not let the debrief eat the last eight minutes --- the homework start is not
optional padding.
\end{multicols}
\end{skillbox}

\boxguard
\begin{skillbox}[AP Practice --- assigned, not scored]{goldbox}
Four AP-style multiple-choice items on parallel boxplots of daily tips for weekday and weekend
shifts --- a context that appears nowhere else in the packet --- then one multi-part free-response
set on the same display. The four MC items cover the four ways this topic is tested: \textbf{1} is
the width-means-count distractor with a fully plausible wrong option; \textbf{2} is a
five-number-summary reading and IQR item; \textbf{3} is mean-versus-median position under right
skew (EK UNC-1.M.2); and \textbf{4} is the item whose correct answer is that \textbf{the boxplot
cannot answer the question} --- whether the weekend distribution is bimodal. That last one is a
real AP item type and the honest limit of the display.

Free response: report the weekday five-number summary, compute the weekend IQR and apply the
$1.5\times$IQR rule to \$158, compare the two shifts in context, and \textbf{part (d) is the
spiral} --- name the display you would ask for instead and say what it would show, which reaches
back to Lessons 1.3 and 1.4.

\textbf{Not scored --- the cover's score column reads \textbf{NA} for this page.} Go over it at
the start of Lesson~1.7 and hold the AP standard out loud: \emph{in context, and comparative, or
it does not count}.
\end{skillbox}

\boxguard
\begin{skillbox}[Homework --- scored, due the first class after two study halls]{goldbox}
Five items in a third set of contexts the lesson has not used --- reading minutes, bus delays,
and pizza deliveries: \textbf{1} build the five-number summary from an ordered list and match it
to one of three pre-drawn boxplots; \textbf{2} the crux head-on, a student claiming the long
right whisker means most values are large; \textbf{3} the $1.5\times$IQR rule with the arithmetic
in a \texttt{work} block; \textbf{4} a full two-group comparison in context from parallel
boxplots; \textbf{5} the AP-style multiple-choice item with a required one-sentence
justification. The packet closes with the Lesson~1.7 preview.

\textbf{Start it in the last eight minutes of the period, in class and alone.} \textbf{Scoring:}
10 points --- 2 per item. \textbf{Item~1 is the diagnostic} and \textbf{item~2 is where the
misconception shows}; point students at 1, 2, and~3 in class, where a wrong start is cheap to
catch, and let 4 and~5 travel home.

\textbf{Paper homework is the default for this course} --- assign this page unless you have
decided otherwise. \textbf{DeltaMath override (occasional):} on the days you assign DeltaMath
instead, it replaces this page, you strike through the score cell on the cover, and this page
stays in the packet as extra practice. Never assign both.
\end{skillbox}

\boxguard
\begin{skillbox}[Watch For (while circulating)]{redbox}
\begin{itemize}[leftmargin=1.2em, nosep]
  \item \textbf{``Most of the players are in the wide part''} (notes~11, 12; AP MC~1; HW~2, HW~5
        --- the target misconception). Probe: \emph{``Count them. How many dots in the wide
        section? Now the narrow one. Say that out loud.''}
  \item \textbf{Estimating the counts instead of counting} (notes~10). The crux only lands if the
        four numbers come out exactly 5, 5, 5, 5. Probe: \emph{``Point at each dot as you say the
        number.''}
  \item \textbf{Reading the box as ``where the data is'' and the whiskers as empty} (notes~5,
        AP MC~1). Probe: \emph{``How many of the twenty are past the right edge of the box?''}
  \item \textbf{Two descriptions parked side by side} (notes~8, 20; HW~4 --- the second
        misconception). Probe: \emph{``I read your sentence about Period 2. Now cross out every
        word that is not comparing.''}
  \item \textbf{Comparison with no context} (notes~20, HW~4): ``the median is higher'' with no
        variable, no units, no groups. Probe: \emph{``Higher what? Higher than what?''}
  \item \textbf{The lone dot read as an error in the data} (notes row~3). It is a real rider with a
        real time. Probe: \emph{``Did somebody mistype it, or did somebody have a bad morning?''}
  \item \textbf{Shape claimed from a boxplot} (notes~14, AP MC~4, HW~5): a boxplot cannot show a
        second peak. Probe: \emph{``Where inside that box would a second hump be drawn?''}
\end{itemize}
Cold-call prompts: \emph{``What fraction of the data is in the left whisker?''} \quad
\emph{``Give me a comparison sentence with a comparative word and a unit in it.''}
\end{skillbox}

\boxguard
\begin{skillbox}[Close \& Start the Homework (8 min)]{goldbox}
\textbf{Students start the homework here, in class, alone and silent --- this is the individual
practice, and the first minutes of it are your best formative read.} Hand the launch first ---
five items on paper, scored, \textbf{due the first class after two study halls} --- and point out
that the AP Practice page is theirs to work and bring to review. Circulate: item~2, the long-right-whisker
claim, is the item that tells you whether the crux row landed. Sort what you see into three piles
as you walk --- corrected with a count (ready); corrected with no number in it (ready, needs the
evidence); agreed with the student (the misconception survived) --- and open Lesson~1.7 by
reteaching the count row if that third pile ran past a third of the room. Close on what changed:
\emph{``You walked in able to read a picture by how big its parts look. You walk out knowing that
in this one the parts are all the same size no matter how they look --- every section holds the
same quarter of the data, and the width is telling you about spread, not about how many.''}
\textbf{Preview:} Lesson~1.7, \emph{The Normal Distribution and $z$-Scores} --- scores from two
completely different tests, and a fair way to say which one was better. Note that everything today
was descriptive: a boxplot describes the sample in front of you and claims nothing about anyone
else.
\end{skillbox}

% ── Teacher notes — one per component, in packet order (four: there is no activity) ──
\begin{teachernote}[Warm-Up]
Five minutes, and item~1 is the one that must be right --- if a student cannot produce $Q_1$ and
$Q_3$ from nine ordered values, row~1 of the notes stalls and everything after it stalls with it.
Circulate during item~1 specifically and fix quartile technique on the spot. \textbf{Debrief
item~1 aloud and leave its five numbers on the board}: the first line of the notes names them,
and pointing at their own numbers while you say \emph{five-number summary} is worth more than the
definition is. Item~2 is a 60-second counting drill; do not let it expand --- it exists only so
that counting is not the obstacle in problem~10. \textbf{Put item~3's four lines of arithmetic on
the board and leave them there for the whole period}: notes row~3's \texttt{work} block runs the
same rule on Route B's 58-minute morning, and the AP Practice free response runs it a third time
on a \$158 shift.
\end{teachernote}

\begin{teachernote}[Guided Notes \& Practice]
\textbf{This one document is the lesson --- pace it 24 / 10, then 8 for the debrief, then 8 for
the homework start.} The notes are a \emph{Main Ideas / Notes} table, not prose: five rows, 20
short problems, and \textbf{in every row you fill the definition lines, write the row's
\texttt{work} block, and work the first problem while the class watches, then the class works the
rest with the pen in their hand.} Rows 1, 2 and~4 are Topic 1.8 and all three read the same
free-throw picture, so the picture is worth 30 seconds of orientation before row~1: twenty dots on
top, the boxplot built from them underneath.

\textbf{Spend the time in One Quarter, the crux row --- it is the last instruction row, straight
before Guided Practice.} Problem~10 is the whole lesson: twenty dots
counted into four sections, coming out 5, 5, 5, 5, even though the sections are 4.5, 5.0, 11.0,
and 23.5 shots wide. \textbf{Every boundary falls inside a gap in the data} (33 to 36, 38 to 41,
49 to 52), so there is no judgment call and no excuse for estimating --- a group that writes
``about 4, about 7'' has lost the discovery, so make them point at each dot. Then
\textbf{problem~11 is the trap} (the classmate who reads width as count) and \textbf{problem~12
is the counterweight} (width reports spread). Take a hand count on 11 before anyone speaks and
pick one \emph{almost}-right answer off a desk --- the debrief opens on it.

Two quieter problems carry more than they look like they do. \textbf{Problem~5} --- how many of
the twenty sit inside the box --- is an early read on the misconception; a student who answers
``most of them'' will fail problem~11. \textbf{Problem~14} --- why the boxplot cannot show the
clump from 30 to 38 --- is the honest limit of the display, and it comes back as the AP Practice
page's MC~4.

\textbf{Two Science Sections is where the pen changes hands.} Problems 15--20, a fresh context,
worked entirely together. \textbf{Problem~20 is the AP-scored form} and the one that cannot be
cut: one paragraph, in context, every claim carrying a comparative word. A paragraph that
describes Period 2 and then describes Period 5 \textbf{earns nothing} on an AP comparison item, no
matter how accurate each half is --- say that out loud before they write, and hold it while they
do. If you only get through 15, 16, and~20, that is the right trade.

\textbf{There is no independent practice set on this page, and that is deliberate --- the homework
is the individual practice.} Guided Practice is the last thing on the page; the period ends with
students starting the homework in front of you, which is where your second formative read comes
from.

\textbf{Debrief (8 min), spoken.} Open on the \emph{almost}-right problem~11, then the count row,
then the comparison standard off the two problem~20 paragraphs you collected. If time is short,
cut the Lesson~1.4 skew callback and the boxplot's-cost item; \textbf{the count row cannot be
cut}. Do not let the debrief eat the last eight minutes.
\end{teachernote}

\begin{teachernote}[AP Practice]
\textbf{This page is not scored} --- the graded practice for this lesson is the homework that
follows it at the back of the packet. Work it as AP-format reps and go over it at the start of
Lesson~1.7. The weekday-versus-weekend tip records appear nowhere else in the packet, so the four
multiple-choice items are the first time students meet this data set.

\textbf{MC item~1 is the one worth reading closely}: option (A) says more weekend shifts fall in
the widest stretch of the display, which is the target misconception dressed as an AP answer
choice. \textbf{MC item~4 is the second one to read}: its correct answer is that the display
cannot tell you whether weekend tips are bimodal, and students who have decided a boxplot ``shows
the shape'' will not even consider it. If either goes badly for more than a handful, spend three
minutes on the counting picture from notes problem~10 before starting 1.7. On the free response,
part~(c) is the highest-value part --- Skill 2.D in the exact form the exam asks it, scored on
comparative language rather than on arithmetic --- and part~(d) is the spiral back to Lessons 1.3
and 1.4.
\end{teachernote}

\begin{teachernote}[Homework]
\textbf{Scored, 10 points (2 per item), due the first class after two study halls.} Launch it in
the last eight minutes so students hit items 1--3 while you can still answer.

\textbf{Item~1 is the diagnostic}: a student who matches the wrong plot has a quartile problem,
not a boxplot problem, and needs two minutes of Lesson 1.5 rather than this lesson repeated.
\textbf{Item~2 is where the misconception shows} --- the long right whisker holds 25 of the 100
employees, exactly as many as the short left one --- and a correction with no number in it is
only half an answer. Item~3's arithmetic is in a \texttt{work} block authored byte-identically in
the blank and the key; require the boundary to be written out, not just the verdict.
\textbf{Score item~4 against the comparison standard}: no comparative word, no credit, however
accurate the two descriptions are. If more than a third of the class lands in that pile, reopen
problem~20's paragraph in the Lesson~1.7 warm-up debrief rather than letting it ride into
inference. Item~5's justification sentence is where you find out whether the crux survived the
drive home.

\textbf{Paper is the default.} This page is what gets assigned unless you decide otherwise on a
given day. On the occasions you assign DeltaMath in its place, strike the score cell on the cover
and tell students this page is now optional practice. Do not assign both.
\end{teachernote}
\end{document}
